\documentclass[10pt, letterpaper]{article}

% Packages:
\usepackage[
  ignoreheadfoot, % set margins without considering header and footer
  top=1 cm, % seperation between body and page edge from the top
  bottom=1 cm, % seperation between body and page edge from the bottom
  left=1 cm, % seperation between body and page edge from the left
  right=1 cm, % seperation between body and page edge from the right
  footskip=1.0 cm, % seperation between body and footer
  % showframe % for debugging 
]{geometry} % for adjusting page geometry
\usepackage{titlesec} % for customizing section titles
\usepackage{tabularx} % for making tables with fixed width columns
\usepackage{array} % tabularx requires this
\usepackage[dvipsnames]{xcolor} % for coloring text
\definecolor{primaryColor}{RGB}{0, 79, 144} % define primary color
\usepackage{enumitem} % for customizing lists
\usepackage{fontawesome5} % for using icons
\usepackage{amsmath} % for math
\usepackage[
  pdftitle={CV - korlando},
  pdfauthor={Kevin Bosirany ORLANDO},
  pdfcreator={LaTeX with RenderCV},
  colorlinks=true,
  urlcolor=primaryColor
]{hyperref} % for links, metadata and bookmarks
\usepackage[pscoord]{eso-pic} % for floating text on the page
\usepackage{calc} % for calculating lengths
\usepackage{bookmark} % for bookmarks
\usepackage{lastpage} % for getting the total number of pages
\usepackage{changepage} % for one column entries (adjustwidth environment)
\usepackage{paracol} % for two and three column entries
\usepackage{ifthen} % for conditional statements
\usepackage{needspace} % for avoiding page brake right after the section title
\usepackage{iftex} % check if engine is pdflatex, xetex or luatex
\usepackage{graphicx}
\usepackage{tikz}

% Ensure that generate pdf is machine readable/ATS parsable:
\ifPDFTeX
  \input{glyphtounicode}
  \pdfgentounicode=1
  % \usepackage[T1]{fontenc} % this breaks sb2nov
  \usepackage[utf8]{inputenc}
  \usepackage{lmodern}
\fi

% Some settings:
\AtBeginEnvironment{adjustwidth}{\partopsep0pt} % remove space before adjustwidth environment
\pagestyle{empty} % no header or footer
\setcounter{secnumdepth}{0} % no section numbering
\setlength{\parindent}{0pt} % no indentation
\setlength{\topskip}{0pt} % no top skip
\setlength{\columnsep}{0cm} % set column seperation
\makeatletter
\let\ps@customFooterStyle\ps@plain % Copy the plain style to customFooterStyle
\patchcmd{\ps@customFooterStyle}{\thepage}{
%  \color{gray}\textit{\small Kevin Bosirany ORLANDO - Page \thepage{} of \pageref*{LastPage}}
}{}{} % replace number by desired string
\makeatother
\pagestyle{customFooterStyle}

\titleformat{\section}{\needspace{4\baselineskip}\bfseries\large}{}{0pt}{}[\vspace{1pt}\titlerule]

\titlespacing{\section}{
  % left space:
  -1pt
}{
  % top space:
  0.3 cm
}{
  % bottom space:
  0.2 cm
} % section title spacing

\renewcommand\labelitemi{$\circ$} % custom bullet points
\newenvironment{highlights}{
  \begin{itemize}[
    topsep=0.10 cm,
    parsep=0.10 cm,
    partopsep=0pt,
    itemsep=0pt,
    leftmargin=0.4 cm + 10pt
  ]
}{
  \end{itemize}
} % new environment for highlights

\newenvironment{highlightsforbulletentries}{
  \begin{itemize}[
    topsep=0.10 cm,
    parsep=0.10 cm,
    partopsep=0pt,
    itemsep=0pt,
    leftmargin=10pt
  ]
}{
  \end{itemize}
} % new environment for highlights for bullet entries


\newenvironment{onecolentry}{
  \begin{adjustwidth}{
    0.2 cm + 0.00001 cm
  }{
    0.2 cm + 0.00001 cm
  }
}{
  \end{adjustwidth}
} % new environment for one column entries

\newenvironment{twocolentry}[2][]{
  \onecolentry
  \def\secondColumn{#2}
  \setcolumnwidth{\fill, 4.6 cm}
  \begin{paracol}{2}
}{
  \switchcolumn \raggedleft \secondColumn
  \end{paracol}
  \endonecolentry
} % new environment for two column entries

\newenvironment{header}{
  \setlength{\topsep}{0pt}\par\kern\topsep\centering\linespread{1.5}
}{
  \par\kern\topsep
} % new environment for the header

\newcommand{\placelastupdatedtext}{% \placetextbox{<horizontal pos>}{<vertical pos>}{<stuff>}
  \AddToShipoutPictureFG*{% Add <stuff> to current page foreground
  \put(
    \LenToUnit{\paperwidth-2 cm-0.2 cm+0.05cm},
    \LenToUnit{\paperheight-1.0 cm}
  ){\vtop{{\null}\makebox[0pt][c]{
    %\small\color{gray}\textit{Dernière mise à jour de Avril 2025}\hspace{\widthof{Last updated in September 2024}}
  }}}%
  }%
}%

% save the original href command in a new command:
\let\hrefWithoutArrow\href

% new command for external links:
\renewcommand{\href}[2]{\hrefWithoutArrow{#1}{\ifthenelse{\equal{#2}{}}{ }{#2 }\raisebox{.15ex}{\footnotesize \faExternalLink*}}}

%%%%%%%%%%%%%%%%%%%%%%%%%%%%%%%%%%%%%%%%%%%%%%%%%%%%%%%%%%%%%%%%%%%%%%%%%%%%%%%%

\begin{document}

\newcommand{\AND}{\unskip
  \cleaders\copy\ANDbox\hskip\wd\ANDbox
  \ignorespaces
}
\newsavebox\ANDbox
\sbox\ANDbox{}

\placelastupdatedtext
\begin{header}
  \vspace{0.3 cm}
  
  \begin{minipage}[c]{0.15\textwidth}
    \centering
    \begin{tikzpicture}
      \clip (0,0) circle (1.5cm);
      \node at (0,0) {\includegraphics[width=3cm]{../../images/avatar.jpg}};
    \end{tikzpicture}
  \end{minipage}%
  \hfill
  \begin{minipage}[c]{0.80\textwidth}
    \textbf{\fontsize{20 pt}{20 pt}\selectfont Kevin Bosirany ORLANDO}

    \vspace{0.3 cm}

    \normalsize
    \mbox{{\color{black}\footnotesize\faMapMarker*}\hspace*{0.13cm}Montpellier, France}%
    \kern 0.25 cm%
    \AND%
    \kern 0.25 cm%
    \mbox{\hrefWithoutArrow{mailto:kevinbosirany@gmail.com}{\color{black}{\footnotesize\faEnvelope[regular]}\hspace*{0.13cm}kevinbosirany@gmail.com}}%
    \kern 0.25 cm%
    \AND%
    \kern 0.25 cm%
    \mbox{\hrefWithoutArrow{tel:+33-762-204-28-40}{\color{black}{\footnotesize\faPhone*}\hspace*{0.13cm}+33 (0)7 62 04 28 40}}%
    \kern 0.25 cm%
    \AND%
    \kern 0.25 cm%
    \mbox{
      \hrefWithoutArrow{https://linkedin.com/in/kbosirany}{
        \color{black}{\footnotesize\faLinkedinIn}\hspace*{0.13cm}kbosirany
      }
    }%
    \kern 0.25 cm%
    \AND%
    \kern 0.25 cm%
    \mbox{\hrefWithoutArrow{https://gitlab.com/kbosirany}{\color{black}{\footnotesize\faGitlab}\hspace*{0.13cm}kbosirany}}%
    \kern 0.25 cm%
    \AND%
    \kern 0.25 cm%
    \mbox{
      \hrefWithoutArrow{https://orcid.org/0009-0009-2784-3108}{
        \color{black}{\footnotesize\faOrcid}\hspace*{0.13cm}0009-0009-2784-3108
      }
    }%
  \end{minipage}
\end{header}

\vspace{0.2 cm}

% Profil

\section{Profil}

\begin{onecolentry}
  Ingénieur en modélisation hydrologique et agronomique avec deux (02) ans d'
  expérience en modélisation et étude d'impact du changement climatique sur les
  ressources en eau.
  Compétences en modélisation des systèmes hydrologiques complexes, en
  développement de modèle sous R et en traitement de données.
  En quête d'un nouveau défi au sein d'un collectif où valoriser et développer
  mes compétences dans un cadre opérationnel.
\end{onecolentry}

% Expériences professionnelles

\section{Expériences professionnelles}

  \begin{twocolentry}
    {
      \textit{Montpellier, France}
      
      \textit{Sept 2023 – Aujourd'hui}
    } {
      \textbf{Ingénieur d'études}

      Chargé de modélisation hydrologique et agronomique

      \textit{INRAE - UMR G-EAU / UMR Innovation}
    }
  \end{twocolentry}

  \vspace{0.10 cm}

  \begin{onecolentry}
    \begin{highlights}
      \item{
        Etude d'impact du changement climatique sur les volumes prélevables du
        bassin versant de l'Aude (Projet Talanoa Water)
        \begin{itemize}
          \item Modélisation hydrologique semi-distribuée du réseau hydrographique
          et hydraulique du bassin versant (GR4J + CemaNeige)
          \item Acquisition et traitement des données climatiques (SAFRAN et DRIAS
          Explore 2), prélèvements en eau (AERMC, DDTM, etc.) et de gestion des
          infrastructures hydrauliques (EDF, VNF, BRL, etc.)
          \item Calibration du modèle et naturalisation des débits
          \item Calcul des volumes prélevables
        \end{itemize}
      }
      \item{
        Développement du modèle Optirrig (Génération, analyse et optimisation des
        scénarios d'irrigation)
        \begin{itemize}
          \item Développement et maintenance du package R du modèle (OptirrigCORE)
          \item Développement et implémentation des formalismes du modèle
          (croissance de la plante, bilan hydrique, etc.)
          \item Acquisition et traitement des données expérimentales
          \item Calibration du modèle pour différentes cultures (maïs, blé, etc.)
        \end{itemize}
      }
      \item Création d'une base de données spatiale des besoins en eau des
      principales cultures irriguées en France sur la période 2000-2100 (Projet
      MAA - PTGE)
    \end{highlights}
  \end{onecolentry}

\vspace{0.2 cm}

% Stages

\section{Stages}

  \begin{twocolentry} {
    \textit{Montpellier, France}
    
    \textit{Mars – Août 2023}
  } {
    \textbf{Stage de fin d'études}
    
    \textit{INRAE - UMR G-EAU / UMR Innovation}
  }
  \end{twocolentry}

  \vspace{0.10 cm}

  \begin{onecolentry}
    \begin{highlights}
      \item Modélisation de la gestion de la ressource en eau du bassin versant de
      l'Aude dans un contexte de changement climatique (Projet Talanoa Water)
    \end{highlights}
  \end{onecolentry}

  \vspace{0.1 cm}

  \begin{twocolentry}
    {
      \textit{Rabat, Maroc}

      \textit{Fév - Juil 2022}
    }
    {
      \textbf{Stage de fin d'études}

      \textit{Institut Agronomique et Vétérinaire Hassan II - Département
      Génie Rural}
    }
  \end{twocolentry}

  \vspace{0.10 cm}

  \begin{onecolentry}
    \begin{highlights}
      \item Détermination de la consommation et de la productivité de l'eau
      dans le bassin hydraulique de l'Oum Er Rbia
    \end{highlights}
  \end{onecolentry}

  \vspace{0.1 cm}

  \begin{twocolentry}
    {
      \textit{Rabat, Maroc}

      \textit{Juil - Sept 2021}
    }
    {
      \textbf{Stage d'insertion en milieu professionnel}

      \textit{Compagnie d'Aménagement et de Développement Industriel (ADI)}
    }
  \end{twocolentry}

  \vspace{0.10 cm}

  \begin{onecolentry}
    \begin{highlights}
      \item Dimensionnement du réseau projeté pour la reconversion en irrigation
      localisée du secteur RD1 du périmètre de la Tessaout aval
    \end{highlights}
  \end{onecolentry}

\vspace{0.2 cm}
  
% Parcours académiques

\section{Parcours académiques}

  \begin{twocolentry}
    {
      \textit{Montpellier, France}

      \textit{Sept 2022 – Sept 2023}
    }
    {
      \textbf{Master 2 Science de l'Eau - Eau et Agriculture}

      \textit{Université de Montpellier - Institut Agro Montpellier - AgroParisTech}
    }
  \end{twocolentry}

  \vspace{0.1 cm}

  \begin{twocolentry}
    {
      \textit{Rabat, Maroc}

      \textit{Sept 2019 – Juil 2022}
    }
    {
      \textbf{Ingénieur Génie Rural - Eau, Environnement et Infrastructures}

      \textit{Institut Agronomique et Vétérinaire Hassan II}
    }
  \end{twocolentry}

  \vspace{0.1 cm}

  \begin{twocolentry}
    {
      \textit{Antananarivo, Madagascar}

      \textit{Sept 2015 – Juil 2019}
    }
    {
      \textbf{Licence Génie Civil - Hydraulique et Aménagement}

      \textit{Ecole Supérieure Polytechnique Antananarivo}
    }
  \end{twocolentry}

\vspace{0.2 cm}

% Compétences

\section{Compétences}
  \begin{onecolentry}
    \textbf{OS:} Windows, Linux
  \end{onecolentry}

  \vspace{0.1 cm}

  \begin{onecolentry}
    \textbf{Bureautique:} Microsoft Office, LibreOffice, etc.
  \end{onecolentry}

  \vspace{0.1 cm}

  \begin{onecolentry}
    \textbf{Programmation:} R, Python (base), Fortran (base)
  \end{onecolentry}

  \vspace{0.1 cm}

  \begin{onecolentry}
    \textbf{Versioning:} Git, Gitlab, Github
  \end{onecolentry}

  \vspace{0.1 cm}

  \begin{onecolentry}
    \textbf{CI/CD:} Gitlab CI, Github Actions
  \end{onecolentry}

  \vspace{0.1 cm}

  \begin{onecolentry}
    \textbf{Autres:} CAO/DAO (AutoCAD), Epanet, etc.
  \end{onecolentry}

\end{document}
